\% !TEX root = ../../buch.tex \% !TEX encoding = UTF-8 \%
\textbackslash section\{Anwendung und Erweiterung der
Padé-Approximation\} \textbackslash label\{pade:section:anwendung\}
\textbackslash kopfrechts\{Anwendung und Erweiterung\}

\textbackslash subsection\{Anwendung der Padé-{[}1/1{]}-Approximation\}
\textbackslash label\{pade:subsection:beispiel\}

Nachdem im vorherigen Abschnitt die Padé-{[}1/1{]}-Approximation
hergeleitet wurde, soll sie nun anhand eines einfachen Zahlenbeispiels
angewendet werden. Ziel ist es, die Genauigkeit der Padé-Approximation
mit jener des entsprechenden Taylor-Polynoms zu vergleichen.

Als Beispiel dient wiederum die Exponentialfunktion
\textbackslash(f(x)=e\^{}x\textbackslash). Der Funktionswert soll für
\textbackslash(x=0.5\textbackslash) approximiert werden.

Der exakte Wert beträgt \textbackslash begin\{equation\} e\^{}\{0.5\} =
1.64872. \textbackslash label\{pade:eq:exp05\}
\textbackslash end\{equation\} Zunächst wird das Taylor-Polynom zweiter
Ordnung verwendet, \textbackslash begin\{equation\} T\_2(x) =
1+x+\textbackslash frac\{x\^{}2\}\{2\}.
\textbackslash label\{pade:eq:taylor2\} \textbackslash end\{equation\}
Durch Einsetzen von \textbackslash(x=0.5\textbackslash) erhält man
\textbackslash begin\{align\} T\_2(0.5) \&= 1 + 0.5 +
\textbackslash frac\{0.5\^{}2\}\{2\} \textbackslash\textbackslash{} \&=
1 + 0.5 + 0.125 \textbackslash\textbackslash{} \&= 1.625.
\textbackslash end\{align\}

Nun wird die Padé-{[}1/1{]}-Approximation verwendet,
\textbackslash begin\{equation\} R\_\{1,1\}(x) =
\textbackslash frac\{1+\textbackslash frac\{x\}\{2\}\}
\{1-\textbackslash frac\{x\}\{2\}\}.
\textbackslash label\{pade:eq:pade11\_beispiel\}
\textbackslash end\{equation\} Für denselben Wert ergibt sich
\textbackslash begin\{align\} R\_\{1,1\}(0.5) \&=
\textbackslash frac\{1+0.25\} \{1-0.25\} \textbackslash\textbackslash{}
\&= \textbackslash frac\{1.25\}\{0.75\} \textbackslash\textbackslash{}
\&= 1.66667. \textbackslash end\{align\}

Zur objektiven Beurteilung beider Verfahren werden nun der absolute und
der relative Fehler berechnet.

Der absolute Fehler ist definiert durch \textbackslash begin\{equation\}
\textbackslash varepsilon\_\{\textbackslash mathrm\{abs\}\} =
\textbar x\_\{\textbackslash mathrm\{exakt\}\}-x\_\{\textbackslash mathrm\{approx\}\}\textbar.
\textbackslash label\{pade:eq:absolut\} \textbackslash end\{equation\}
Der relative Fehler ergibt sich zu \textbackslash begin\{equation\}
\textbackslash varepsilon\_\{\textbackslash mathrm\{rel\}\} =
\textbackslash frac\{\textbar x\_\{\textbackslash mathrm\{exakt\}\}-x\_\{\textbackslash mathrm\{approx\}\}\textbar\}
\{\textbar x\_\{\textbackslash mathrm\{exakt\}\}\textbar\}.
\textbackslash label\{pade:eq:relativ\} \textbackslash end\{equation\}

Die Ergebnisse sind in
Tabelle\textasciitilde\textbackslash ref\{pade:tab:vergleich\}
zusammengefasst.

\textbackslash begin\{table\}{[}ht{]} \textbackslash centering
\textbackslash caption\{Vergleich der Approximationen für
\textbackslash(x=0.5\textbackslash).\}
\textbackslash label\{pade:tab:vergleich\}
\textbackslash begin\{tabular\}\{lccc\} \textbackslash hline Verfahren
\& Approximation \& Absoluter Fehler \& Relativer
Fehler\textbackslash\textbackslash{} \textbackslash hline Exakter Wert
\& 1.64872 \& 0 \& 0\textbackslash\textbackslash{} Taylor
\textbackslash(T\_2\textbackslash) \& 1.62500 \& 0.02372 \&
1.44\textbackslash,\textbackslash\%\textbackslash\textbackslash{} Padé
\textbackslash({[}1/1{]}\textbackslash) \& 1.66667 \& 0.01795 \&
1.09\textbackslash,\textbackslash\%\textbackslash\textbackslash{}
\textbackslash hline \textbackslash end\{tabular\}
\textbackslash end\{table\}

Bereits dieses einfache Beispiel zeigt, dass beide Verfahren mit wenigen
Rechenschritten eine gute Näherung liefern. Die
Padé-{[}1/1{]}-Approximation weist jedoch sowohl einen kleineren
absoluten als auch einen kleineren relativen Fehler auf als das
Taylor-Polynom gleicher Ordnung.

Dieses Ergebnis verdeutlicht den grundlegenden Vorteil der
Padé-Approximation. Beide Verfahren verwenden ausschliesslich dieselben
Informationen der Taylor-Reihe. Während die Taylor-Approximation diese
Informationen direkt zu einem Polynom zusammenfasst, nutzt die
Padé-Approximation sie zur Konstruktion einer rationalen Funktion.
Dadurch kann häufig eine höhere Genauigkeit erreicht werden.

Ein einzelnes Zahlenbeispiel erlaubt jedoch nur eine begrenzte Aussage
über die Qualität einer Approximation. Daher wird im folgenden Abschnitt
das Verhalten beider Verfahren über einen grösseren Wertebereich
miteinander verglichen.

\textbackslash subsection\{Vergleich der Approximationen\}
\textbackslash label\{pade:subsection:vergleich\}

Das vorherige Zahlenbeispiel zeigt den Unterschied zwischen Taylor- und
Padé-Approximation lediglich für einen einzelnen Funktionswert. Für die
Beurteilung eines Approximationsverfahrens ist jedoch insbesondere das
Verhalten über einen grösseren Wertebereich von Interesse.

Abbildung\textasciitilde\textbackslash ref\{pade:fig:vergleich\} zeigt
den Verlauf der Exponentialfunktion, des Taylor-Polynoms zweiter Ordnung
sowie der Padé-{[}1/1{]}-Approximation.

\%Bild 2, Teil3, (vergleich), Taylor vs Pade 1-1

In der unmittelbaren Umgebung des Entwicklungspunktes stimmen beide
Approximationen nahezu mit der Exponentialfunktion überein. Mit
zunehmendem Abstand wird jedoch deutlich, dass die Padé-Approximation
den Verlauf der Exponentialfunktion besser beschreibt als das
Taylor-Polynom.

Der Grund hierfür liegt in der mathematischen Struktur der
Approximation. Während ein Taylor-Polynom ausschliesslich aus Potenzen
der Variablen besteht, besitzt die Padé-Approximation zusätzlich ein
Nennerpolynom. Dadurch stehen weitere Freiheitsgrade zur Verfügung,
welche eine bessere Anpassung an den Funktionsverlauf ermöglichen.

Die Padé-{[}1/1{]}-Approximation besitzt allerdings auch
Einschränkungen. Da ihr Nenner
\textbackslash(1-\textbackslash frac\{x\}\{2\}\textbackslash) für
\textbackslash(x=2\textbackslash) verschwindet, besitzt die
Approximation an dieser Stelle eine Polstelle. Die Exponentialfunktion
selbst besitzt keine Polstellen. Man spricht deshalb von einer
künstlichen Polstelle, welche ausschliesslich durch die Approximation
entsteht.

Dieses Beispiel zeigt, dass auch Padé-Approximationen nicht exakt sind.
Dennoch liefern sie in vielen Anwendungen eine deutlich bessere Näherung
als Taylor-Polynome gleicher Ordnung. Eine weitere Verbesserung kann
durch Padé-Approximationen höherer Ordnung erreicht werden.

\textbackslash subsection\{Padé-Approximationen höherer Ordnung\}
\textbackslash label\{pade:subsection:hoehere\}

Die bisher betrachtete Padé-{[}1/1{]}-Approximation verwendet lediglich
die ersten drei Koeffizienten der Taylor-Reihe. Soll die Genauigkeit
weiter verbessert werden, müssen zusätzliche Informationen der
Taylor-Reihe berücksichtigt werden. Dies geschieht durch
Padé-Approximationen höherer Ordnung.

Allgemein besitzt eine Padé-Approximation vom Typ
\textbackslash({[}m/n{]}\textbackslash) die Form
\textbackslash begin\{equation\} R\_\{m,n\}(x) = \textbackslash frac
\{a\_0+a\_1x+\textbackslash cdots+a\_mx\^{}m\}
\{1+b\_1x+\textbackslash cdots+b\_nx\^{}n\}.
\textbackslash label\{pade:eq:mn\} \textbackslash end\{equation\} Der
Zähler besitzt den Grad \textbackslash(m\textbackslash), der Nenner den
Grad \textbackslash(n\textbackslash). Da der konstante Term des
Nennerpolynoms auf \textbackslash(1\textbackslash) normiert wird,
enthält die Approximation insgesamt \textbackslash(m+n+1\textbackslash)
unbekannte Koeffizienten.

Zur Bestimmung dieser Koeffizienten wird wiederum die Taylor-Reihe der
ursprünglichen Funktion verwendet. Allgemein gilt
\textbackslash begin\{equation\} f(x)-R\_\{m,n\}(x) = O(x\^{}\{m+n+1\}).
\textbackslash label\{pade:eq:ordnung2\} \textbackslash end\{equation\}
Dies bedeutet, dass die Taylor-Reihe der Funktion und die
Reihenentwicklung der Padé-Approximation bis einschliesslich des Terms
\textbackslash(x\^{}\{m+n\}\textbackslash) übereinstimmen. Erst ab
höheren Potenzen treten Unterschiede auf.

Als Beispiel wird nun die Padé-Approximation vom Typ
\textbackslash({[}2/2{]}\textbackslash) betrachtet. Sie besitzt die Form
\textbackslash begin\{equation\} R\_\{2,2\}(x) = \textbackslash frac
\{a\_0+a\_1x+a\_2x\^{}2\} \{1+b\_1x+b\_2x\^{}2\}.
\textbackslash label\{pade:eq:pade22\_ansatz\}
\textbackslash end\{equation\}

Im Vergleich zur Padé-{[}1/1{]}-Approximation treten nun insgesamt fünf
unbekannte Koeffizienten auf. Nach der Entwicklung des Nenners, dem
Ausmultiplizieren sowie dem Koeffizientenvergleich erhält man das
entsprechende lineare Gleichungssystem. Die vollständige Herleitung
erfolgt analog zur Padé-{[}1/1{]}-Approximation und wird deshalb an
dieser Stelle nicht vollständig ausgeschrieben.

Für die Exponentialfunktion ergeben sich die Koeffizienten
\textbackslash begin\{align\} a\_0 \&= 1,\textbackslash\textbackslash{}
a\_1 \&= \textbackslash frac12,\textbackslash\textbackslash{} a\_2 \&=
\textbackslash frac1\{12\},\textbackslash\textbackslash{} b\_1 \&=
-\textbackslash frac12,\textbackslash\textbackslash{} b\_2 \&=
\textbackslash frac1\{12\}. \textbackslash end\{align\}

Damit erhält man die Padé-{[}2/2{]}-Approximation
\textbackslash begin\{equation\} R\_\{2,2\}(x) = \textbackslash frac
\{1+\textbackslash frac\{x\}\{2\}+\textbackslash frac\{x\^{}2\}\{12\}\}
\{1-\textbackslash frac\{x\}\{2\}+\textbackslash frac\{x\^{}2\}\{12\}\}.
\textbackslash label\{pade:eq:pade22\} \textbackslash end\{equation\}

Im Vergleich zur Padé-{[}1/1{]}-Approximation werden nun mehr
Koeffizienten der Taylor-Reihe exakt reproduziert. Dadurch verbessert
sich die Genauigkeit der Approximation über einen grösseren Bereich.

Abbildung\textasciitilde\textbackslash ref\{pade:fig:pade22\} zeigt den
Vergleich zwischen der Exponentialfunktion sowie den
Padé-Approximationen erster und zweiter Ordnung.

\%Bild 3, Teil3, (pade22), e\^{}x vs Pade 1-1 u 2-2

Wie zu erwarten, nähert sich die Padé-{[}2/2{]}-Approximation dem
exakten Funktionsverlauf über einen deutlich grösseren Bereich an als
die Padé-{[}1/1{]}-Approximation. Mit steigender Ordnung verbessert sich
die Approximation im Allgemeinen weiter, gleichzeitig nimmt jedoch auch
der Rechenaufwand zur Bestimmung der Koeffizienten zu.

Im letzten Abschnitt dieses Kapitels werden die wichtigsten
Anwendungsgebiete der Padé-Approximation vorgestellt sowie ihre Vor- und
Nachteile zusammengefasst.

\textbackslash subsection\{Anwendungen der Padé-Approximation\}
\textbackslash label\{pade:subsection:anwendungen\}

Nachdem die mathematischen Grundlagen der Padé-Approximation vorgestellt
wurden, stellt sich die Frage nach ihrer praktischen Bedeutung. Obwohl
das Verfahren ursprünglich im Rahmen der Approximationstheorie
entwickelt wurde, findet es heute in zahlreichen wissenschaftlichen und
technischen Anwendungen Verwendung. Überall dort, wo Funktionen
effizient und gleichzeitig mit hoher Genauigkeit ausgewertet werden
müssen, kann die Padé-Approximation ihre Stärken ausspielen.

Eine wichtige Anwendung findet sich in der numerischen Mathematik. Viele
elementare Funktionen wie die Exponentialfunktion, der natürliche
Logarithmus oder trigonometrische Funktionen werden in numerischen
Algorithmen nicht direkt berechnet, sondern durch Approximationen
ersetzt. Padé-Approximationen liefern dabei häufig genauere Ergebnisse
als Taylor-Polynome gleicher Ordnung und werden deshalb in verschiedenen
numerischen Verfahren eingesetzt.

Ein weiteres bedeutendes Anwendungsgebiet ist die Regelungstechnik.
Viele technische Systeme besitzen sogenannte Totzeiten. Zwischen einer
Änderung der Eingangsgrösse und der Reaktion des Systems vergeht dabei
eine bestimmte Zeit. Im Laplace-Bereich wird eine Totzeit durch
\textbackslash begin\{equation\} e\^{}\{-sT\}
\textbackslash label\{pade:eq:totzeit\} \textbackslash end\{equation\}
beschrieben.

Da diese Exponentialfunktion analytische Berechnungen erschwert, wird
sie häufig durch eine Padé-Approximation ersetzt. Für eine
Padé-{[}1/1{]}-Approximation ergibt sich
\textbackslash begin\{equation\} e\^{}\{-sT\} \textbackslash approx
\textbackslash frac\{1-\textbackslash frac\{sT\}\{2\}\}
\{1+\textbackslash frac\{sT\}\{2\}\}.
\textbackslash label\{pade:eq:totzeit\_pade\}
\textbackslash end\{equation\} Dadurch entsteht eine rationale
Übertragungsfunktion, welche sich mit den Methoden der klassischen
Regelungstechnik wesentlich einfacher untersuchen lässt.

Abbildung\textasciitilde\textbackslash ref\{pade:fig:totzeit\} zeigt
beispielhaft den Vergleich zwischen der exakten Sprungantwort einer
Totzeit und der entsprechenden Padé-{[}1/1{]}-Approximation.

\%Bild 4, Teil3, (totzeit), Pade von Totzeit

Auch bei der numerischen Lösung gewöhnlicher und partieller
Differentialgleichungen werden Padé-Approximationen eingesetzt. Häufig
wird zunächst eine Taylor-Reihe der Lösung bestimmt und anschliessend in
eine Padé-Approximation umgewandelt. Dadurch kann die Lösung oftmals
über einen grösseren Bereich mit höherer Genauigkeit beschrieben werden.

Darüber hinaus finden Padé-Approximationen zahlreiche Anwendungen in der
theoretischen Physik. Insbesondere in der Quantenmechanik, der
statistischen Physik und der Strömungsmechanik treten
Reihenentwicklungen auf, deren Konvergenzbereich begrenzt ist. Durch die
Umwandlung dieser Reihen in rationale Funktionen lassen sich häufig
genauere numerische Ergebnisse erzielen.

Ein weiterer Vorteil der Padé-Approximation zeigt sich darin, dass sie
nicht auf reelle Argumente beschränkt ist. Da sie für Funktionen mit
konvergenter Potenzreihe konstruiert werden kann, besitzen insbesondere
holomorphe Funktionen eine Padé-Approximation. Damit liefert die
Padé-Approximation nicht nur für reelle, sondern auch für komplexe Werte
eine Näherung.

Ein anschauliches Beispiel liefert erneut die Exponentialfunktion. Für
ein komplexes Argument \textbackslash(z=ix\textbackslash) gilt
\textbackslash(e\^{}\{ix\}=\textbackslash cos(x)+i\textbackslash sin(x)\textbackslash).
Eine Padé-Approximation der Exponentialfunktion kann deshalb auch zur
Approximation von Sinus und Kosinus verwendet werden.
Abbildung\textasciitilde\textbackslash ref\{pade:fig:komplex\} zeigt den
Vergleich zwischen den exakten Funktionen und den aus der
Padé-{[}1/1{]}-Approximation gewonnenen Real- und Imaginärteilen.

\%Bild 5, Teil3, (komplex), Pade Komplex

Ein weiterer Vorteil besteht darin, dass die Padé-Approximation die aus
der Taylor-Reihe bekannte Ableitungsinformation erhält. Da die
Reihenentwicklung der Padé-Approximation bis zur geforderten Ordnung mit
jener der ursprünglichen Funktion übereinstimmt, stimmen auch die
entsprechenden Ableitungen am Entwicklungspunkt überein. Die Ableitung
einer Padé-Approximation liefert somit wiederum eine Approximation der
Ableitung der ursprünglichen Funktion.

Bei herkömmlichen Interpolationsverfahren ist diese Eigenschaft nicht
automatisch gegeben. Soll eine Funktion und gleichzeitig ihre Ableitung
an bestimmten Stellen approximiert werden, muss beispielsweise eine
spezielle Methode wie die Hermite-Interpolation verwendet werden.

\textbackslash subsection\{Vor- und Nachteile\}
\textbackslash label\{pade:subsection:vorteile\}

Wie jedes Approximationsverfahren besitzt auch die Padé-Approximation
sowohl Vorteile als auch Einschränkungen.

Ein wesentlicher Vorteil besteht darin, dass ausschliesslich die
Koeffizienten der Taylor-Reihe benötigt werden. Zusätzliche
Informationen über die Funktion sind nicht erforderlich. Durch die
Darstellung als rationale Funktion kann der Verlauf vieler Funktionen
genauer beschrieben werden als mit einem Taylor-Polynom gleicher
Ordnung.

Insbesondere ausserhalb der unmittelbaren Umgebung des
Entwicklungspunktes liefern Padé-Approximationen häufig deutlich
kleinere Approximationsfehler. Darüber hinaus können rationale
Funktionen Eigenschaften beschreiben, welche mit Polynomen grundsätzlich
nicht dargestellt werden können, beispielsweise asymptotisches
Verhalten.

Demgegenüber stehen jedoch auch einige Nachteile. Je höher die Ordnung
der Padé-Approximation gewählt wird, desto grösser wird das zu lösende
lineare Gleichungssystem und damit auch der Rechenaufwand. Zudem können
künstliche Polstellen entstehen, welche in der ursprünglichen Funktion
nicht vorhanden sind.

Die wichtigsten Eigenschaften beider Verfahren sind in
Tabelle\textasciitilde\textbackslash ref\{pade:tab:vor\_nachteile\}
zusammengefasst.

\textbackslash begin\{table\}{[}ht{]} \textbackslash centering
\textbackslash caption\{Vergleich zwischen Taylor- und
Padé-Approximation.\} \textbackslash label\{pade:tab:vor\_nachteile\}
\textbackslash begin\{tabular\}\{lll\} \textbackslash hline Eigenschaft
\& Taylor \& Padé\textbackslash\textbackslash{} \textbackslash hline
Grundlage \& Taylor-Reihe \& Taylor-Reihe\textbackslash\textbackslash{}
Darstellung \& Polynom \& rationale
Funktion\textbackslash\textbackslash{} Genauigkeit \& lokal sehr gut \&
häufig auf grösserem Bereich\textbackslash\textbackslash{} Polstellen \&
keine \& künstliche Polstellen möglich\textbackslash\textbackslash{}
Rechenaufwand \& gering \& höher\textbackslash\textbackslash{}
Anwendungen \& lokale Approximation \& Numerik, Regelungstechnik,
Physik\textbackslash\textbackslash{} \textbackslash hline
\textbackslash end\{tabular\} \textbackslash end\{table\}
